\documentclass[letterpaper,11pt]{article}

\usepackage[utf8]{inputenc}
\usepackage[T1]{fontenc}
\usepackage{lmodern}
\usepackage[spanish,es-tabla]{babel}
\usepackage[margin=2.3cm]{geometry}
\usepackage{setspace}
\usepackage{graphicx}
\usepackage{booktabs}
\usepackage{tabularx}
\usepackage{array}
\usepackage{xcolor}
\usepackage{enumitem}
\usepackage{amsmath}
\usepackage{fancyhdr}
\usepackage{microtype}
\usepackage[hidelinks]{hyperref}

% -------------------- FORMATO --------------------
\setstretch{1.08}
\setlength{\parindent}{1em}
\setlength{\parskip}{4pt}
\setlength{\headheight}{15pt}

\pagestyle{fancy}
\fancyhf{}
\fancyfoot[C]{\thepage}
\renewcommand{\headrulewidth}{0.4pt}
\fancyhead[L]{\small Equipo 57}
\fancyhead[R]{\small Equidad lingüística en tecnologías de voz}


\definecolor{azul}{RGB}{28,65,105}
\definecolor{gris}{RGB}{245,246,248}

% -------------------- COMANDOS --------------------
\newcommand{\metric}[3]{%
\begin{minipage}[t]{0.235\linewidth}
\colorbox{gris}{%
\parbox[t]{0.91\linewidth}{%
\centering
\vspace{3pt}
{\Large\bfseries\color{azul}#1}\par
{\small\bfseries #2}\par
{\scriptsize #3}
\vspace{3pt}
}}
\end{minipage}
}

\begin{document}

% ============================================================
% PORTADA / ENCABEZADO
% ============================================================

\begin{center}

{\Large\bfseries\color{azul}
HACIA LA EQUIDAD LINGÜÍSTICA EN TECNOLOGÍAS DE VOZ
}

\vspace{0.2cm}

{\large
Benchmark multimodal de reconocimiento automático del habla
}

\vspace{0.15cm}

{\normalsize
Resumen ejecutivo para CLEAR
}

\end{center}

\vspace{0.3cm}

\noindent
\colorbox{gris}{
\parbox{\dimexpr\linewidth-2\fboxsep\relax}{
\small
\textbf{Equipo:} 57 -- Proyecto Integrador, Grupo 10\\
\textbf{Integrantes:} Alejandro Agustín Ahedo González;
Andrés Roberto Osuna González;
Sandra Luz Cervantes Espinoza\\
\textbf{Organización receptora:} CLEAR
(Coalition for Language Equity in AI Research)\\
\textbf{Asesoría:} Dra. Liliana Ibeth Barbosa Santillán\\
\textbf{Fecha:} 11 de noviembre de 2026
}}

% ============================================================
\section*{\color{azul}1. La decisión ejecutiva}

La equidad lingüística en inteligencia artificial exige que
las personas puedan interactuar con estas tecnologías en sus
idiomas sin que la calidad del servicio dependa injustificadamente
de cuánto está representada su lengua en los datos digitales.

Cuando un sistema de reconocimiento automático del habla
funciona mejor para unos idiomas que para otros, reproduce
una barrera de acceso: algunas comunidades reciben
transcripciones más útiles, mientras otras enfrentan más
errores, menor confiabilidad y una experiencia tecnológica
desigual. Los promedios globales pueden ocultar esas diferencias.

CLEAR (Coalition for Language Equity
in AI Research) sitúa la equidad entre idiomas como un
problema que debe medirse y abordarse con evidencia.
En su análisis sobre la desigualdad de tokenización,
CLEAR muestra cómo determinadas decisiones técnicas
pueden imponer costos desiguales a distintas lenguas [1].

Ese análisis se centra en tokenizadores de modelos de lenguaje, 
buscando que los usuarios de IA puedan
hacer visibles las disparidades lingüísticas y aplicar
una evaluación transparente al reconocimiento automático
del habla (ASR).

El proyecto construye un flujo reproducible
para comparar sistemas ASR, medir sus errores por idioma
y aportar evidencia para seleccionar y mejorar tecnologías
de voz con un criterio explícito de equidad lingüística. Utiliza el corpus abierto FLEURS, diseñado como benchmark
multilingüe para investigación en tecnologías del habla [2],
un pipeline de normalización acústica y textual y métricas
complementarias WER y CER.

El objetivo es
identificar qué idiomas están mejor o peor atendidos bajo
condiciones comparables y qué aspectos requieren más datos,
evaluación o adaptación.

% ============================================================
\section*{\color{azul}2. Evidencia técnica disponible}

En la evaluación final, la configuración seleccionada fue
\texttt{openai/whisper-large-v3} con decodificación base
(greedy), perteneciente a la familia de modelos de
reconocimiento multilingüe de voz descrita por
Radford et al. [3].

Sobre el conjunto de prueba reservado de ocho idiomas,
el notebook reporta los siguientes resultados:

\vspace{0.2cm}

\noindent
\metric{4.91\%}{CER en prueba}{Ocho idiomas}
\hfill
\metric{8.03\%}{WER en prueba}{Idiomas con espacios}
\hfill
\metric{33.75\%}{Coincidencia exacta}{Texto normalizado}
\hfill
\metric{1.491 s}{Latencia promedio}{Por audio}

\vspace{0.3cm}

La búsqueda Beam Search con cinco haces obtuvo mejores
cifras en validación (CER 2.96\% y WER 4.01\%), pero la
selección final priorizó el equilibrio observado entre
error y tiempo en prueba, no el resultado de validación
por sí solo.


\textbf{Implicación para la equidad lingüística.}

Estos resultados constituyen una primera línea de evidencia,
no una demostración de equidad entre idiomas.

El conjunto final cubre ocho idiomas y el tamaño de prueba
es reducido. Por ello, no permite afirmar que el sistema
trate de manera justa a todas las comunidades lingüísticas
ni extrapolar el desempeño a lenguas no evaluadas. Sí demuestra 
por qué es necesario medir y reportar los
errores por idioma y elegir métricas compatibles con cada
sistema de escritura.

El CER permite comparar sistemas de escritura sin espacios;
el WER debe reservarse para idiomas en los que la
segmentación por palabras sea válida.

En el avance exploratorio, el uso de CER corrigió la lectura
del japonés: Whisper pasó de un WER artificialmente elevado
de 74.99\% a un CER de 6.21\%. Esto ilustra por qué la métrica
debe adaptarse a la lengua. Investigaciones sobre evaluación
de ASR en japonés también documentan dificultades relacionadas
con segmentación y variantes ortográficas [5].

Los resultados respaldan un piloto controlado de evaluación
y transcripción asistida, pero no justifican todavía un
despliegue productivo sin validación adicional.

% ============================================================
\section*{\color{azul}3. Contribución a la equidad lingüística
y valor para CLEAR}

La contribución central para CLEAR es aportar un instrumento
reproducible para observar desigualdades de desempeño en
tecnologías de voz y orientar decisiones de investigación,
evaluación y mejora.

La equidad no se presume por el carácter multilingüe de un
modelo: debe examinarse con resultados desagregados,
criterios adecuados a cada idioma y participación de personas
con conocimiento lingüístico y cultural de las comunidades
evaluadas.

El benchmark puede ayudar a detectar brechas y priorizar
nuevas evaluaciones; por sí solo no resuelve la falta de datos,
representación o recursos para cada lengua.

El pipeline implementado contempla:

\begin{itemize}[leftmargin=*,itemsep=2pt]
    \item Estandarización del audio a 16 kHz, canal mono
    y formato \texttt{float32}.
    \item Normalización Unicode y reglas de procesamiento
    textual.
    \item Selección reproducible de muestras mediante
    semilla fija.
    \item Evaluación con métricas WER y CER.
    \item Uso de caché persistente para reutilizar datos
    procesados.
\end{itemize}

En la preparación de datos se documentó que la caché
persistente redujo el tiempo de preparación de 106 minutos
a menos de un minuto en ejecuciones posteriores.
Esta mejora corresponde a la preparación del experimento
y no debe interpretarse como ahorro mensual de operación
productiva.

\subsection*{Beneficios cuantificables}

Para dimensionar el posible beneficio operativo se modela
un escenario ilustrativo de 10,000 clips mensuales.

Bajo los supuestos del Anexo A, un minuto de trabajo manual
por clip, 1.491 segundos de inferencia por clip y revisión
humana del 15\% de los resultados a un minuto por clip
revisado, el tiempo de trabajo estimado baja de 166.67
a 29.14 horas mensuales.

El ahorro potencial es de:

\begin{center}
{\Large\bfseries\color{azul}
137.53 horas mensuales (82.5\%)
}
\end{center}

Es una proyección para dimensionar un piloto, no un ahorro
observado ni un compromiso de capacidad. Deberá recalibrarse
con tiempos reales, volumen y tasa de revisión del sponsor.

\subsection*{Beneficios intangibles}

La solución fortalece la toma de decisiones al ofrecer
métricas desagregadas por idioma en lugar de depender
exclusivamente de promedios globales.

Mejora la calidad del proceso al incorporar normalización
y criterios de evaluación adecuados a cada sistema de escritura.

Simplifica el trabajo de investigación mediante pipelines,
caché y resultados reproducibles.

Finalmente, aporta una base para convertir la equidad
lingüística en un criterio técnico verificable: no basta
con que un sistema admita varios idiomas, es necesario
comprobar cómo funciona en cada uno y transparentar dónde falla.

% ============================================================
\section*{\color{azul}4. Inversión, operación y condiciones
de implementación}

Se recomienda avanzar por etapas. Primero, cerrar la
validación técnica del benchmark, después, ejecutar un
piloto acotado con volúmenes, idiomas y criterios de
aceptación acordados; y solo entonces decidir una
integración productiva.

El prototipo se ejecutó en Google Colab con GPU Tesla T4
(14.56 GB de VRAM reportados). El notebook recomienda
una GPU de al menos 16 GB, por lo que la infraestructura
observada no debe asumirse como dimensionamiento definitivo
para producción.

Como estimación preliminar de planeación (no cotización
de proveedor) se considera:

\begin{itemize}[leftmargin=*,itemsep=2pt]
    \item 120 horas de desarrollo a MXN \$600 por hora:
    MXN \$72,000.
    \item 40 horas de implementación a MXN \$600 por hora:
    MXN \$24,000.
    \item 30 horas-GPU mensuales a MXN \$45 por hora:
    MXN \$1,350.
    \item Almacenamiento, monitoreo y respaldos:
    MXN \$650 mensuales.
\end{itemize}

El costo operativo estimado es de MXN \$2,000 mensuales.
El costo total de propiedad estimado para 12 meses es
de MXN \$120,000.

\vspace{0.2cm}

\begin{center}
\begin{tabular}{lr}
\toprule
\textbf{Concepto} & \textbf{Importe MXN}\\
\midrule
Desarrollo & \$72,000\\
Implementación & \$24,000\\
Operación mensual & \$2,000\\
TCO estimado a 12 meses & \$120,000\\
\bottomrule
\end{tabular}
\end{center}

Esta estimación no incluye impuestos, servicios de API,
integración empresarial, soporte extendido, anotación
adicional ni revisión humana. Los supuestos deben
sustituirse por cotizaciones y tarifas reales antes
de autorizar presupuesto.

% ============================================================
\section*{\color{azul}5. Riesgos principales y mitigación}

\small
\begin{tabularx}{\linewidth}{
    @{}
    >{\raggedright\arraybackslash}p{0.28\linewidth}
    X
    @{}
}
\toprule
\textbf{Riesgo} & \textbf{Medida de mitigación propuesta}\\
\midrule

Cobertura lingüística limitada &
El experimento final reporta ocho idiomas.
Ampliar progresivamente el conjunto y validar con hablantes
o revisores competentes antes de extrapolar resultados
a otras lenguas o dialectos.\\

Errores de inferencia y reproducibilidad &
Los registros del notebook muestran fallos en dos audios
(alemán e italiano) asociados a entradas que activan el modo
de audio largo. Corregir el manejo de duración, reejecutar
la prueba completa y conservar logs, versiones y hashes
de datos.\\

Sesgo de medición &
Mantener normalización específica por idioma; usar CER
para escrituras sin espacios y WER solo donde la
segmentación léxica sea válida. Reportar métricas por idioma,
no únicamente promedios.\\

Capacidad e infraestructura &
Dimensionar GPU y concurrencia mediante pruebas de carga.\\

Costos, privacidad y licencias &
Establecer presupuesto por volumen, monitorear consumo,
limitar datos a fuentes autorizadas y revisar licencias
y condiciones de uso antes de incorporar datos o
servicios externos.\\

\bottomrule
\end{tabularx}
\normalsize

% ============================================================
\section*{\color{azul}6. Transferencia y continuidad}

Al sponsor se propone transferir:

\begin{enumerate}[leftmargin=*]
    \item Repositorio de código y notebooks reproducibles.
    \item Pipeline de ingesta, normalización y evaluación.
    \item Resultados WER/CER y desglose por idioma.
    \item Documentación de configuración, dependencias,
    limitaciones y criterios de selección.
    \item Archivos de resultados y visualizaciones.
    \item Guía de ejecución y mantenimiento.
\end{enumerate}

El repositorio reportado por el equipo es
\href{https://github.com/Luz026/VoiceBenchmark}{VoiceBenchmark}.
El código se identifica con licencia MIT y el corpus
FLEURS con CC-BY-4.0, sujeto a conservar las atribuciones
y verificar las condiciones aplicables.

La recomendación ejecutiva es autorizar un piloto técnico condicionado con enfoque de
equidad lingüística, no un despliegue general inmediato.

Sus condiciones de salida deben ser:

\begin{itemize}[leftmargin=*]
    \item Ejecución completa sin muestras fallidas.
    \item Resultados estables en una prueba ampliada.
    \item Validación lingüística por idioma.
    \item Medición de costos reales.
    \item Aprobación de privacidad y licenciamiento.
\end{itemize}

Con estas condiciones, CLEAR podrá utilizar el benchmark
como mecanismo de observación y seguimiento de brechas
en sistemas ASR, aportando evidencia para que la innovación
en voz avance hacia una participación más equitativa de
las comunidades lingüísticas.

% ============================================================
% ANEXO A
% ============================================================

\clearpage
\section*{\color{azul}Anexo A. Supuestos y cálculo de
beneficios y costos}

\subsection*{A.1 Beneficio mensual de ahorro de tiempo}

El siguiente cálculo es un escenario para
dimensionar el piloto.

\textbf{Supuestos:}

\begin{itemize}[leftmargin=*]
    \item Volumen mensual: 10,000 clips.
    \item Trabajo manual de referencia: 1 minuto por clip.
    \item Inferencia automatizada: 1.491 segundos por clip.
    \item Control humano: revisión del 15\% de los clips,
    a 1 minuto por clip revisado.
\end{itemize}

\textbf{Cálculo:}

\begin{align*}
T_{\text{manual}}
&=
\frac{10{,}000 \times 1\ \mathrm{min}}{60}
=
166.67\ \mathrm{h/mes}
\\
T_{\text{inferencia}}
&=
\frac{10{,}000 \times 1.491\ \mathrm{s}}{3{,}600}
=
4.14\ \mathrm{h/mes}
\\
T_{\text{revision}}
&=
\frac{10{,}000 \times 0.15 \times 1\ \mathrm{min}}{60}
=
25.00\ \mathrm{h/mes}
\\
T_{\text{automatizado}}
&=
4.14 + 25.00
=
29.14\ \mathrm{h/mes}
\\
Ahorro
&=
166.67 - 29.14
=
137.53\ \mathrm{h/mes}
\\
Reduccion
&=
\frac{137.53}{166.67}\times 100
=
82.5\%
\end{align*}

El ahorro de tiempo no equivale automáticamente a
ahorro de caja: representa capacidad laboral potencial
liberada. Para monetizarlo se requiere conocer el costo
laboral real, la proporción de tiempo efectivamente
reasignable y si se evita gasto incremental.

No se reporta reducción de errores mensual ni incremento
de ventas porque los avances no contienen una línea base
operativa mensual que permita calcularlos de forma defendible.

% -------------------- COSTOS --------------------

\subsection*{A.2 Estimación de inversión y costo total
de propiedad}

\begin{center}
\small
\begin{tabularx}{\linewidth}{@{}Xrr@{}}
\toprule
\textbf{Concepto} &
\textbf{Base supuesta} &
\textbf{Importe MXN}\\
\midrule
Desarrollo & 120 h $\times$ \$600/h & \$72,000\\
Implementación & 40 h $\times$ \$600/h & \$24,000\\
GPU mensual & 30 h $\times$ \$45/h & \$1,350\\
Almacenamiento, monitoreo y respaldos &
Cuota mensual supuesta & \$650\\
Operación mensual & GPU + servicios básicos & \$2,000\\
Operación 12 meses & 12 $\times$ \$2,000 & \$24,000\\
\midrule
\textbf{TCO estimado a 12 meses} &
\$72,000 + \$24,000 + \$24,000 &
\textbf{\$120,000}\\
\bottomrule
\end{tabularx}
\end{center}

Las tarifas son supuestos de planeación, no precios
observados ni cotizaciones. No incluyen impuestos,
servicios de API, integración empresarial, soporte
de terceros, anotación de nuevos corpus, seguridad
especializada ni costo de revisores humanos.

El costo de GPU debe recalcularse con el proveedor,
tipo de acelerador, concurrencia, almacenamiento
y volumen real.

% -------------------- ACEPTACION --------------------

\subsection*{A.3 Condiciones de aceptación recomendadas
para el piloto}

\begin{enumerate}[leftmargin=*]
    \item Reejecutar la evaluación completa sin errores
    de procesamiento y confirmar el denominador efectivo
    de las métricas.

    \item Congelar versiones de modelo, librerías, datos
    y reglas de normalización.

    \item Publicar WER y CER por idioma, además de promedios
    macro y micro claramente identificados; incluir
    dispersión y tamaño de muestra.

    \item Validar muestras con hablantes competentes y,
    cuando sea posible, representantes de las comunidades.

    \item Medir latencia, VRAM, costo por mil audios y tasa
    de revisión en la infraestructura objetivo.
\end{enumerate}

% ============================================================
% ANEXO B: REFERENCIAS
% ============================================================

\clearpage
\section*{\color{azul}Anexo B. Declaración de uso de IA
y referencias}

Los avances técnicos del proyecto documentan el uso de
herramientas de IA como apoyo para estructuración,
depuración de código, evaluación, métricas y redacción.

El equipo conserva la responsabilidad de revisar,
ejecutar y validar los resultados, el código y las
conclusiones antes de su difusión. Esta declaración
deberá cotejarse con el registro final de herramientas
y contribuciones del equipo antes de la publicación
institucional.

\subsection*{Referencias bibliográficas}

\begin{enumerate}[leftmargin=*,itemsep=7pt]

\item CLEAR -- Coalition for Language Equity in AI Research.
(2026). \textit{The Tokenizer Tax}. Portal oficial de CLEAR.
Estudio sobre disparidades de tokenización entre idiomas.
Disponible en:
\url{https://clear-languages.org/blog/the-tokenizer-tax}

\item Conneau, A., Ma, M., Khanuja, S., Zhang, Y.,
Axelrod, V., Dalmia, S., Riesa, J., Rivera, C., y
Bapna, A. (2022).
\textit{FLEURS: Few-shot Learning Evaluation of Universal
Representations of Speech}.
arXiv:2205.12446.
\url{https://arxiv.org/abs/2205.12446}

\item Radford, A., Kim, J. W., Xu, T., Brockman, G.,
McLeavey, C., y Sutskever, I. (2023).
\textit{Robust Speech Recognition via Large-Scale Weak
Supervision}.
Proceedings of the 40th International Conference on
Machine Learning, 202, 28492--28518.
\url{https://proceedings.mlr.press/v202/radford23a.html}

\item Thennal, D. K., James, J., Gopinath, D. P.,
y Ashraf, K. (2025).
\textit{Advocating Character Error Rate for Multilingual
ASR Evaluation}.
Findings of the Association for Computational Linguistics:
NAACL 2025, 4941--4950.
\url{https://doi.org/10.18653/v1/2025.findings-naacl.277}

\item Karita, S., Sproat, R., y Ishikawa, H. (2023).
\textit{Lenient Evaluation of Japanese Speech Recognition:
Modeling Naturally Occurring Spelling Inconsistency}.
Proceedings of the Workshop on Computation and Written
Language (CAWL 2023), 61--70.
\url{https://doi.org/10.18653/v1/2023.cawl-1.8}

\item Ahia, O., Kumar, S., Gonen, H., Kasai, J.,
Mortensen, D. R., Smith, N. A., y Tsvetkov, Y. (2023).
\textit{Do All Languages Cost the Same? Tokenization in
the Era of Commercial Language Models}.
Proceedings of the 2023 Conference on Empirical Methods
in Natural Language Processing, 9904--9923.
\url{https://doi.org/10.18653/v1/2023.emnlp-main.614}

\end{enumerate}

% ============================================================
% FIN DEL DOCUMENTO
% ============================================================

\end{document}
